\subsection{An exact component trace after restriction of scalars and a Tate twist}

\begin{proposition}[P-F72: a literal geometric trace for each quadratic source value]
\label{prop:nonstandardtatetrace}
Retain the actual source values $\alpha_\pm(q)$ and the polynomials
$A,B,E$ of P-F69.  For every $q=3^e$, $e\ge1$, put
\[
 K=\mathbb Q(\sqrt{A(q)}),\qquad d=117,\qquad Q=q^d,
 \qquad
 \gamma=\frac B2\bigl(E+q^2(1+q^8)\sqrt{A(q)}\bigr)\in K.
                                                               \tag{P-F72a}
\]
There is a simple ordinary abelian surface
$\mathcal A/\mathbb F_Q$ with $K$ acting by rational
$\mathbb F_Q$-endomorphisms, and an abelian variety
\[
          \mathcal R=\operatorname{Res}_{\mathbb F_Q/\mathbb F_q}
                          \mathcal A,\qquad \dim\mathcal R=234,
                                                               \tag{P-F72b}
\]
with the following property.  Rational algebraic correspondences
over $\mathbb F_q$ induce a Frobenius-equivariant idempotent on
$H^{117}(\mathcal R_{\overline{\mathbb F}_q},\mathbb Q_\ell)$,
for every $\ell\ne3$.  Its image $W_\ell$ has a commuting
$K$-action and is free of rank $2^{117}$ over
$K_\ell=K\otimes\mathbb Q_\ell$.  Moreover,
\[
 \operatorname{Tr}_{K_\ell}
       \bigl(\operatorname{Frob}_q\mid W_\ell(26)\bigr)
       =q^{-26}\gamma=\alpha(q),                              \tag{P-F72c}
\]
where $\alpha(q)$ is the element of $K$ whose two real
conjugates are $\alpha_+(q),\alpha_-(q)$.
Consequently, after extending coefficients and choosing either
embedding of $K$, the corresponding source value is the literal
geometric Frobenius trace on the associated rank-$2^{117}$
summand.  The summand $W_\ell(26)$ is pure of weight $65$.

Here $(26)$ is the ordinary integral Tate twist over
$\mathbb F_q$.  There is no additional rational scalar in
(P-F72c).  The projector is asserted to be idempotent on
cohomology; no unstated idempotence in a Chow ring is needed.
\end{proposition}

\begin{proof}
\par\medskip\noindent\textbf{1. An integral real trace with an odd extension degree.}
The integer $q$ is odd, so the ten-term sum
$1+q+\cdots+q^9$ is even.  Thus $B/2$ is an integer.
The radical $\sqrt{A(q)}$ is an algebraic integer, and
(P-F72a) therefore defines an algebraic integer $\gamma$.
Its two conjugates satisfy
\[
       \gamma_\pm=\tfrac12\beta_\pm
                 =q^{26}\alpha_\pm(q),\qquad
       0<\gamma_\pm<243q^{52}<q^{58}<2\sqrt Q.                \tag{P-F72d}
\]
The bound is the explicit estimate in (P-F69e), with
$243<q^6$ for $q\ge3$; it does not use the old choice of $Q$.
P-F66 proves $A(q)>0$ and $A(q)\equiv2\pmod3$, so $K$ is
real quadratic.  The coefficient of $\sqrt{A(q)}$ in
$\gamma$ is nonzero, and hence $\mathbb Q(\gamma)=K$.
At every prime of $K$ above $3$ one has
\[
                         \gamma\equiv1.                     \tag{P-F72e}
\]
Indeed $B\equiv1$, $E\equiv2$, and the radical term is
divisible by $q^2$; division by $2$ is legitimate modulo
such a prime.

Choose a root $\pi$ of $X^2-\gamma X+Q$.
At both real embeddings this quadratic has negative
discriminant by (P-F72d), with conjugate roots of absolute
value $\sqrt Q$.  Thus $\pi$ is an integral $Q$-Weil number.
Since $\gamma=\pi+Q/\pi$ generates $K$, the field
$L=\mathbb Q(\pi)=K(\pi)$ is quartic CM.
The Honda--Tate construction and the local-invariant check
in P-F69 apply with this $\gamma$ and this $Q$: at every
prime above $3$, Hensel's lemma gives one unit root and
one root of valuation $v(Q)$ by (P-F72e).  Every local
division-algebra invariant is zero, there are no real
places, and therefore
\[
 \operatorname{End}_{\mathbb F_Q}^0(\mathcal A)=L,
 \qquad \dim\mathcal A=2,
 \qquad
 \det_{K_\ell}(X-\operatorname{Frob}_Q\mid H^1)=
                         X^2-\gamma X+Q.                    \tag{P-F72f}
\]
The surface is ordinary because $\gamma$ is a unit above
the characteristic.  Precisely the primary inputs used
here are Tate, Bourbaki~352, Theorem~1(i)--(ii), equations
(1)--(2) and Remark~2, printed p.~96, and Waterhouse,
Proposition~7.1, printed p.~553, as cited and checked in
P-F69.  In particular, no assertion about absolute
simplicity or base extension is being imported.

\par\medskip\noindent\textbf{2. The all-$H^1$ summand of an actual Weil restriction.}
The Weil restriction in (P-F72b) exists and is an abelian
variety.  Over $\overline{\mathbb F}_q$ it is the product
of the $d$ conjugate surfaces $\mathcal A_i$, indexed by
the embeddings of $\mathbb F_Q$ into the algebraic closure:
\[
             \mathcal R_{\overline{\mathbb F}_q}
                 \simeq\prod_{i=0}^{d-1}\mathcal A_i.
                                                               \tag{P-F72g}
\]
For existence, compatibility with base change, and this
product description we use Milne, \emph{On the arithmetic
of abelian varieties}, \S1, printed pp.~177--178:
\url{https://www.jmilne.org/math/articles/1972a.pdf}.
The finite extension is separable and the surface is
projective, so the stated hypotheses apply.

Write $V_i=H^1(\mathcal A_i,\mathbb Q_\ell)$.
The exterior-algebra description of abelian-variety
cohomology and its Galois compatibility give the
K\"unneth summand
\[
       U_\ell=\bigotimes_{i=0}^{d-1}V_i
              \ \subseteq\ H^d(\mathcal R_{\overline{\mathbb F}_q},
                                      \mathbb Q_\ell).        \tag{P-F72h}
\]
We use Milne, \emph{Abelian varieties}, Theorem~15.1(a)--(b),
printed p.~27, and Remark~15.5, printed p.~29, in the
author's corrected original article:
\url{https://www.jmilne.org/math/xnotes/AVs.pdf}.
In particular $[2]^*$ acts as $2^j$ on $H^j$ of each
surface.  On its $i$th factor the rational polynomial
\[
             p_i=\prod_{\substack{0\le j\le4\\j\ne1}}
                    \frac{[2]_i^*-2^j}{2-2^j}
                                                               \tag{P-F72i}
\]
is therefore the cohomological projector onto degree one.
The commuting product $p=\prod_i p_i$ projects precisely
onto (P-F72h), and is zero on every other multidegree.
Each operation is induced by a rational algebraic
correspondence: scalar multiples of the diagonal and
graphs of multiplication maps suffice.

\par\medskip\noindent\textbf{3. The equal-embedding projector and its descent.}
Transport the fixed embedding $K\hookrightarrow L$ to
each conjugate surface.  Geometric Frobenius permutes these
copies cyclically; its $d$th power commutes with every
copy because those endomorphisms are defined over
$\mathbb F_Q$.  Let $a=A(q)$ and let $u_i$ denote the
action of $\sqrt a$ on $V_i$, and hence on the $i$th
tensor factor of $U_\ell$.  On $U_\ell$ these operators
commute and satisfy $u_i^2=a$.  Put
\[
             e_{\rm eq}=
                 \prod_{i=1}^{d-1}
                       \frac{1+u_0u_i/a}{2}.                 \tag{P-F72j}
\]
Each factor is an idempotent.  Over a field splitting $K$,
the tensor factors are labeled by signs
$\varepsilon_i\in\{+1,-1\}$, the eigenvalue of $u_i$
being $\varepsilon_i\sqrt a$.  Formula (P-F72j) equals
one exactly when all $\varepsilon_i$ agree and equals
zero otherwise.  It is consequently invariant under
every permutation of the factors, despite the asymmetric
choice of index $0$ in the formula.  Define
\[
                           W_\ell=e_{\rm eq}U_\ell.
                                                               \tag{P-F72k}
\]
There are two surviving sign summands, each of dimension
$2^d$ by (P-F72f).  On $W_\ell$ all the $u_i$ coincide.
Their common action defines a $K_\ell$-module, free of
rank $2^d$, and it commutes with geometric Frobenius.
Equivalently, the natural quotient from the tensor product
over $\mathbb Q_\ell$ to $\bigotimes_{K_\ell,i}V_i$
restricts to an isomorphism on $W_\ell$: after splitting
$K_\ell$, both sides are exactly the two all-equal-sign
summands.  This specifies the tensor product over the
coefficient field used in the next step.

This projector and coefficient action are geometric
in the stated correspondence sense.  If a transported
quasi-endomorphism $\sqrt a$ has denominator $m$, use
$m^{-1}[m\sqrt a]^*$ on its factor and sandwich with
$p$.  On the all-$H^1$ summand it gives exactly $u_i$;
no assertion of additivity of pullback in other
cohomological degrees is required.  Thus $pe_{\rm eq}p$
is represented by a rational correspondence over
$\mathbb F_Q$.  Each of its $d$ Galois conjugates acts
as this same projector: $p$ is invariant, and the
equal-sign condition in (P-F72j) is invariant.
Push the correspondence forward along the finite
base-change map to $\mathcal R\times\mathcal R$ over
$\mathbb F_q$ and divide by $d$.
After base change its action is the average of the
$d$ conjugate actions, and is hence precisely
$pe_{\rm eq}p$.  This proves the claimed descent and
cohomological idempotence.  The same averaging applied
to $u_0pe_{\rm eq}p$ gives the common $K$-action over
$\mathbb F_q$, since $u_i=u_0$ on the image.

\par\medskip\noindent\textbf{4. The cyclic trace and the Tate twist.}
We spell out the sign and coefficient convention.
After choosing identifications of successive $K_\ell$-modules
by geometric Frobenius transport, the action on the
equal-embedding tensor is the cyclic operator
\[
 C(v_0\otimes\cdots\otimes v_{d-1})
       =Fv_{d-1}\otimes v_0\otimes\cdots\otimes v_{d-2},
       \qquad F=\operatorname{Frob}_Q\mid H^1(\mathcal A),
                                                               \tag{P-F72l}
\]
multiplied by $(-1)^{d-1}$ from reordering $d$ factors
of cohomological degree one in the cup product.  Reversing
the indexing of the cycle merely moves the displayed
$F$ to another tensor position.  Its product around
the cycle is always geometric $\operatorname{Frob}_Q$,
not its inverse.  Because $d=117$ is odd, the sign is $+1$.

For a free module with basis $v_1,v_2$ over either field
factor of $K_\ell$, the diagonal matrix coefficients
of $C$ vanish unless all the tensor-basis indices agree.
Their sum is $F_{11}+F_{22}$.  Therefore, over $K_\ell$,
\[
                       \operatorname{Tr}(C)=\operatorname{Tr}(F)
                                              =\gamma.        \tag{P-F72m}
\]
This calculation also proves the identity when
$K_\ell$ is a product of fields, by applying it on
each factor.  The geometric Frobenius action on
$\mathbb Q_\ell(1)$ is multiplication by $q^{-1}$:
the arithmetic action on the compatible roots of
unity is the $q$th power, and geometric Frobenius is
its inverse.  Tensoring (P-F72m) with the ordinary
Tate line $\mathbb Q_\ell(26)$ now proves (P-F72c).

Finally, the $d$th power of geometric Frobenius on
$W_\ell$ is the tensor product of the $d$ conjugate
$Q$-Frobenius actions.  Every eigenvalue of each
factor has complex absolute value $Q^{1/2}$ by
(P-F72f).  Hence every eigenvalue of that $d$th power
has absolute value $Q^{d/2}=q^{d^2/2}$, and every
Frobenius eigenvalue on $W_\ell$ has absolute value
$q^{d/2}$.  The twist reduces the weight by $52$,
giving $d-52=65$, as asserted.
\end{proof}

\paragraph{What this exact trace construction does and does not identify.}
The construction uses a different Honda--Tate surface
from P-F69, over $\mathbb F_{q^{117}}$, and then a
$234$-dimensional Weil restriction over $\mathbb F_q$.
It is not a relabeling of the rational scalar in
(P-F69c).  The full rational trace of the twisted
summand is $\alpha_+(q)+\alpha_-(q)$; either individual
value is a trace after choosing one embedding of the
quadratic coefficient field.  The projector has
rational coefficients and the component rank is
$2^{117}$.  Neither effectiveness of the projector
nor a small-dimensional model is asserted.

All choices are fiberwise in $q=3^e$; quadratic
base-extension compatibility is ruled out below.
No family over a parameter line,
no equality between source multiplication and the
tensor-induced correspondences, and no identification
of a source word with Frobenius are asserted.
In particular P-F66's obstruction to making the
quadratic value itself a pure Frobenius eigenvalue
is untouched: (P-F72c) is a trace of many eigenvalues,
all of one weight, rather than a claim about the
weight of their sum.  The remaining source-product
and theta-operation comparisons retain their separate
quantifiers.

\begin{corollary}[The exact trace models still fail base-extension compatibility]
For every $e\ge1$, the twisted projected cohomology
$W_\ell(26)$ constructed at $q=3^e$, after quadratic
base extension, is not isomorphic as a Frobenius
representation to the corresponding construction at
$q^2=3^{2e}$.  This holds for every choice of the
Honda--Tate surfaces and the stated projectors.
\end{corollary}

\begin{proof}
Since $\gcd(2,d)=1$, squaring the cyclic operator
in (P-F72l) leaves a single cycle of tensor positions.
In a tensor basis its diagonal-coefficient contraction
runs once around that cycle and encounters $F$ twice.
The same calculation as (P-F72m) therefore gives
\[
            \operatorname{Tr}_{K_\ell}(C^2)
                   =\operatorname{Tr}_{K_\ell}(F^2)
                   =\gamma^2-2Q.                             \tag{P-F72n}
\]
There is no new sign: the original $d$-cycle already
has sign $+1$ on the all-degree-one summand.
The full rational trace of the twice-applied geometric
Frobenius on the Tate-twisted summand is consequently
\[
       q^{-52}\bigl(\gamma_+^2+\gamma_-^2-4Q\bigr)<0,          \tag{P-F72o}
\]
because (P-F72d) gives
$\gamma_+^2+\gamma_-^2<2q^{116}<2Q$.
On the other hand, the full rational trace of the
construction at $q^2$ is
$\alpha_+(q^2)+\alpha_-(q^2)>0$.
Isomorphic Frobenius representations have equal traces,
which proves the assertion.  This comparison concerns
the specified projected and twisted cohomology, not an
unproved assertion about isogenies of the ambient
Weil restrictions.
\end{proof}
